%% kapa: exam=1 round=26 year=2015 period=2 form=A subject=appraisal_law subject_ko=부동산관계법규 seq=26 number=26
%% source: Q-Net 감정평가사 기출문제 (cst00309) artlSeq=5138001 file=extracted/1008/5138001/제26회_감정평가사 1차_기출/제26회_감정평가사 1차_기출/2교시_감평_A형.pdf
%% conversion: kordoc 4.24.51 -> tools/gen1.py (auto); review=proofread-sonnet; flags=none
\begin{question}{26}
국유재산법령상의 내용으로 옳지 않은 것은?
\begin{choices}
  \item 행정재산을 사용허가한 때에는 매년 사용료를 징수한다.
  \item 정부는 정부출자기업체를 새로 설립하려는 경우 일반재산을 현물출자할 수 있다.
  \item 사권이 설정된 재산은 판결에 따라 취득하는 경우에도 그 사권이 소멸된 후가 아니면 국유재산으로 취득하지 못한다.
  \item 일반재산은 대부할 수 있다.
  \item 행정재산의 사용료가 면제되는 경우도 있다.
\end{choices}
\end{question}
\answer{3}
